\section{Proof of Lemma \ref{lem:key}: computation of a representation of the Kleene closure}

Recall that one important component in the decidability of $\GPtwo$ is the computation of a representation of the Kleene closure of a semilinear set.
Recall from Section \ref{sec:prelims} that $\projp{n}{\bfv}$ denotes the projection operator that maps a vector $\bfv$ to its first $n$ entries.
The formal statement was:

\medskip

For every integer linear system $\cQ(\bfx): \bfA \bfx = \bfc$ with $\bfc \neq \bfzero$,
where $\bfA \in \bbZ^{m \times n}$ and $\bfc \in \bbZ^{m}$,
there exists $\bfAtilde \in \bbZ^{(n+t) \times (n+k)}$ with $\norm{\bfAtilde} = 1$
such that
\begin{equation*}
    \kstar{\sol{\cQ(\bfx)}}\ =\ \sol{\cQtilde(\bfx)},
\end{equation*}
where $\cQtilde(\bfx)$ is the following Boolean combination of integer linear systems
\begin{equation*}    
    \exists \bfy_1, \bfy_2\ 
    (\bfx\ =\ \projp{n}{\bfy_1} + \bfy_2)\ \land\ 
    (\bfAtilde\bfy_1\ =\ \bfzero)\ \land\ 
    (\bfA\bfy_2\ =\ \bfzero)\ \land\ 
    ((\bfy_1=\bfzero) \to (\bfy_2=\bfzero)),
\end{equation*}
$k = n\cdot(2D+1)^{m} + 1$,
$t \le 2nk$, and
$D = n \cdot \norm{\bfA} \cdot \left((n+1)\cdot \norm{\bfA} + \norm{\bfc} + 1\right)^{m}$.
Moreover, $\bfAtilde$ can be computed in time $2^\bigO{\log n + \log K + m^2}$,
where $K := \max\left(\norm{\bfA}, \norm{\bfc}\right)$.

\medskip

We prove the lemma by considering \emph{minimal solutions} of the integer linear system.
Recall that $\vle$ means pointwise inequality between vectors.
\begin{definition}[Minimal elements]
    For a set $\cS \subseteq \bbN^n$ and $\bfv \in \cS$,
    we say that $\bfv$ is \emph{minimal} in $\cS$
    if, for every $\bfu \in \cS$ with $\bfu \neq \bfv$,
     $\bfu \not\vle \bfv$.
\end{definition}

Note that the following properties about minimal elements are trivial.
\begin{lemma}\label{lem:min_property}
    Let $\cS$ be a nonempty subset of $\bbN^n$.
    \begin{enumerate}
        \item There exists $\bfv \in \cS$ such that $\bfv$ is minimal in $\cS$.
        \item For every $\bfv \in \cS$, there exists a minimal element $\bfu \in \cS$,
            such that $\bfv - \bfu \in \bbN^n$.
    \end{enumerate}
\end{lemma}

For an integer linear system $\cQ(\bfx)$,
let $\solmin{\cQ(\bfx)}$ denotes the set of minimal solutions,
i.e., $\solmin{\cQ(\bfx)} := \setc{\bfv \in \sol{\cQ(\bfx)}}{\bfv\ \text{is minimal in}\ \sol{\cQ(\bfx)}}$.
Note that the maximal entry of minimal solutions is bounded.

\begin{lemma}[\cite{Pottier}, also see Proposition 4. in~\cite{semilinear}]\label{lem:minimal}
    For every integer linear system $\cQ(\bfx): \bfA \bfx = \bfc$,
    where $\bfA \in \bbZ^{m \times n}$ and $\bfc \in \bbZ^{m}$,
    for every $\bfv \in \solmin{\cQ(\bfx)}$,
     $\norm{\bfv} \le \left((n+1)\cdot\norm{\bfA} + \norm{\bfc} + 1\right)^{m}$.
\end{lemma}

We begin by turning to the automata-theoretic approach for the solutions of integer linear systems.

\begin{definition}[Acyclic, unambiguous, and simple finite automata] \label{def:simpleplus}
As usual, a finite automaton can be viewed as a 
directed graph where the vertices are the states
and the edges are the transitions.

Such an automaton is \emph{acyclic} if the graph does not contain a cycle. 
It is \emph{unambiguous} if every word has 
at most one accepting run.
    We say that a finite automaton is \emph{simple}
    if it is unambiguous, acyclic, and has exactly one accepting state.
\end{definition}

\begin{definition}[Parikh Vector]
    For a word $w = a_{t_1}\cdots a_{t_\ell} \in \Sigma^*$,
    the \emph{Parikh vector} of $w$
    is a vector in $\bbN^{n}$, defined as
    $\parikh{w} := \sum_{i \in \intsinterval{\ell}} \unit{t_i}$.
    Intuitively, the $i^{th}$ entry in $\parikh{w}$ is the number of occurrences of $a_i$ in $w$.

For a finite automaton $\cM$, the \emph{Parikh image} of $\cM$ is defined as
\begin{equation*}
\parikh{\cM}
\ :=\ 
\setc{\parikh{w}}{\text{$\cM$ accepts $w$}}.
\end{equation*}
\end{definition}

We can encode a finite subset of solutions of a linear system $\cQ(\bfx):\bfA \bfx = \bfc$ as the Parikh image of a finite automaton $\cM$.
More precisely, $\parikh{\cM}$ contains all the minimal solutions of $\cQ(\bfx)$.
The main idea behind the construction of $\cM$ is as follows.
Suppose $\cQ(\bfx)$ has a solution
and $\bfA\in \bbZ^{m\times n}$.
By Lemma~\ref{lem:minimal},
the components of every minimal solution are bounded above by  
$L:=\left((n+1)\cdot\norm{\bfA} + \norm{\bfc} + 1\right)^{m}$.
We can construct a finite automaton $\cM$ where the set of states is $\setc{\bfA\bfv}{\bfv \in [0,L]^n}$
and the transitions are
$\langle \bfA\bfv,a_i, \bfA(\bfv+\bfe_i) \rangle$ for every $1\le i \le n$.
The initial state is $\bfzero$ and the accepting state is $\bfc$.
Obviously, $\parikh{\cM}$ contains all the minimal solutions. In fact, it contains all solutions that are bounded by $L$. 
Note also that
$\cM$ can be constructed in exponential time.
With a little more work, we can make $\cM$ simple, a property that will be useful when analyzing the Kleene star of $\sol{\cQ}$.

\begin{lemma}\label{lem:ilp_to_nfa}
For every integer linear system $\cQ(\bfx): \bfA \bfx = \bfc$,
where $\bfA \in \bbZ^{m \times n}$ and $\bfc \in \bbZ^{m}$,
there exists a simple finite automaton $\cM$
with $k$ states and $t$ transitions,
where $k = n\cdot L\cdot(2D+1)^{m} + 1$, 
$t \le 2nk$,
$L = \left((n+1)\cdot\norm{\bfA} + \norm{\bfc} + 1\right)^{m}$
and $D = n \cdot \norm{\bfA} \cdot L$,
    such that
    \begin{equation*}
        \solmin{\cQ(\bfx)}
        \ \subseteq\ 
        \parikh{\cM}
        \ \subseteq\ 
        \sol{\cQ(\bfx)}.
    \end{equation*}
    Moreover, $\cM$ can be constructed in time exponential in $n,m,\norm{\bfA},\norm{\bfc}$.
\end{lemma}

\begin{proof}
If $\bfc = \bfzero$, then the system $\cQ(\bfx)$ admits the
solution $\bfzero$.
The finite automaton that accepts only the empty word is the desired $\cM$.
So we concentrate on the case when $\bfc\neq \bfzero$.

The automaton $\cM$ is very similar to the one described in the intuition preceding the lemma.
We fix a finite alphabet of $n$ symbols $\Sigma := \set{a_1, \ldots, a_{n}}$.
To ensure acyclicity, we modify the automaton $\cM$ to accept only words from $a_1^*a_2^*\cdots a_n^*$. The automaton will also implement a counter that counts the number of occurrences of each $a_i$ up to $L$. Such a counter can be implemented using the states of $\cM$.
We now give the details.

Let $\cD$ be the set $\intinterval{-D}{D}^{m}$.
\begin{itemize}
\item 
The set of states is 
$\{q_{acc}\}\cup(\intsinterval{n}\times\intsinterval{L} \times \cD)$.
\item
The initial state is $(0,0,\bfzero)$.
\item 
For each $0\leq i<j \leq n$, $0\leq t\leq L-1$,
and $\bfv \in \cD$,
we have the transitions:
\begin{itemize}
\item 
if $\bfv + \bfA\unit{i} \in \cD$,
then $(\tuple{i,t, \bfv}, a_i, \tuple{i,t+1, \bfv + \bfA\unit{i}}) \in \Delta$, 
\item 
if $\bfv + \bfA\unit{i}  = \bfc$,
then $(\tuple{i, t,\bfv}, a_i, q_{acc}) \in \Delta$.

\item 
if $\bfv + \bfA\unit{j} \in \cD$,
then $(\tuple{i,t, \bfv}, a_j, \tuple{j,1, \bfv + \bfA\unit{j}}) \in \Delta$, 
\item 
if $\bfv + \bfA\unit{j}  = \bfc$,
then $(\tuple{i, t,\bfv}, a_j, q_{acc}) \in \Delta$.
\end{itemize}
\end{itemize}
Obviously, $\cM$ has only one accepting state.
The automaton $\cM$ is acyclic since it can only go from the state
$\tuple{i,t, \bfv}$ to $\tuple{j,t',\bfu}$
when $i< j$ or when $i=j$ and $t<t'$.
It is also unambiguous since
for every state $\tuple{i,t,\bfv}$ and $a\in \Sigma$, there is at most one state $\tuple{i',t',\bfu}$
such that $\tuple{i,t,\bfv},a,\tuple{i',t',\bfu}$ is a transition.
Finally, it is routine to verify that
$\parikh{\cM}\subseteq \sol{\cQ}$.
That $\solmin{\cQ}\subseteq \parikh{\cM}$
follows from Lemma~\ref{lem:minimal}.
\end{proof}

Next, we consider the Kleene star of sets.
Recall the definition:

\medskip

For every $\cS \subseteq \bbN^n$,
the \emph{Kleene star} of $\cS$ is defined as
\begin{equation*}
    \kstar{\cS} := \setc{\sum_{s \in \cS'} s}{\text{$\cS'$ is a finite multisubset of $\cS$}}.
\end{equation*}

\medskip

In the following we fix a finite alphabet of $n$ symbols $\Sigma:=\{a_1,\ldots,a_n\}$
and assume that every finite automaton is over this alphabet.
We need more terminology concerning finite automata.
Note that for every acyclic finite automaton $\cM$ (Definition \ref{def:simpleplus})
a run $r$ of $\cM$ uses each transition at most once.
Therefore, we can represent $r$ by a Boolean vector with dimension $\abs{\Delta}$.
Recall that for a unambiguous finite automaton $\cM$ and an accepted word $w$, there exists exactly one accepting run of $\cM$ on $w$.
Let $\bbB:=\{0,1\}$.

\begin{definition}
    For every simple finite automaton $\cM$ with $t$ transitions,
    for every accepted word $w$ of $\cM$,
    letting $r$ be the unique accepting run of $\cM$ on $w$,
    the \emph{path vector} of $w$, denoted by $\mpath{w}$,
    is the vector in $\bbB^t$ defined as follows:
    for $1 \le i \le t$, the $i^{th}$ entry of $\mpath{w}$ is 1
    if and only if $r$ contains the $i^{th}$ transition.
\end{definition}

\begin{lemma}\label{lem:path}
For every simple finite automaton $\cM$ with $t$ transitions,
there exists a matrix $\bfB_\cM \in \bbZ^{n \times t}$,
which can be computed in time $\bigO{nt}$,
such that
\begin{itemize}
\item $\norm{\bfB_\cM} = 1$, and
\item for every word $w$ accepted by $\cM$, $\parikh{w} = \bfB_\cM\cdot\mpath{w}$.
\end{itemize}
\end{lemma}

\begin{proof}
    We construct the matrix $\bfB$ as follows.
    For $1 \le i \le n$ and $1 \le j \le t$,
    if the $t^{th}$ transition is of the form $(q, a_i, p)$,
    then $(B_\cM)_{i, j} = 1$.
    Otherwise, $(B_\cM)_{i, j} = 0$.
    It is trivial to verify that
    $\parikh{w} = \bfB\cdot\mpath{w}$.
\end{proof}

We now prove a  lemma regarding the Kleene closure of the Parikh image.
The proof employs a similar idea as in the argument for Theorem~1 in \cite{tree}.

\begin{lemma}\label{lem:kstar_nfa}
    For every simple finite automaton $\cM$ with $k$ states $q_1, \ldots,q_k$ and $t$ transitions,
    there exists a matrix $\bfA_\cM \in \bbZ^{k \times t}$,
    which can be computed in time $\bigO{kt}$,
    such that
    \begin{itemize}
        \item $\norm{\bfA_\cM} = 1$, and
        \item $\kstar{\parikh{\cM}} =
            \sol{\cQ_\cM(\bfx)}$,
            where
            $\cQ_\cM(\bfx) := \exists \bfy\ (\bfx = \projp{n}{\bfy}) \land (\tilde{\bfA}_\cM \bfy = \bfzero)$ and
            $\bfAtilde_\cM = 
            \left[
                \begin{array}{r|r}
                    -\bfI & \bfB_\cM \\
                    \hline
                    \bfzero& \bfA_\cM
                \end{array}
            \right]$.
    \end{itemize}
\end{lemma}

\begin{proof}
    We construct the matrix $\bfA_\cM$ as follows.
    For $1 \le i \le k$ and $1 \le j \le t$,
    if $q_i$ is the initial state or the accepting state, then
    $\left(A_\cM\right)_{i, j} = 0$.
    Otherwise,
    \begin{equation*}
        \left(A_\cM\right)_{i, j}\ :=\ 
        \begin{cases}
            1,  &\text{if the $j^{th}$ transition is of the form $(p, a, q_i)$} \\
            -1, &\text{if the $j^{th}$ transition is of the form $(q_i, a, p)$} \\
            0,  &\text{otherwise}.
        \end{cases}
    \end{equation*}
    Note that since $\cM$ is acyclic, there is no transition of the form $(q_i, a, q_i)$.

    We first claim that for every accepted word $w$, 
    we have
    \begin{equation*}
        \bfA_\cM \cdot \mpath{w} = \bfzero.    
    \end{equation*}
    We will show that $\left(\bfA_\cM \cdot \mpath{w}\right)_i = 0$,    for $1 \le i \le k$.
    If $q_i$ is the initial state or the accepting state,
    by definition, $\left(A_\cM\right)_{i, j} = 0$ for every $1 \le j \le k$,
    hence, $\left(\bfA_\cM \cdot \mpath{w}\right)_i = 0$.

    If $q_i$ does not appear in the accepting run of $w$,
    then $\mpath{w}_j=0$ if
    the $j^{th}$ transition is going into or coming out of state $q_i$.
    Hence, $\left(\bfA_\cM \cdot \mpath{w}\right)_i = 0$.

    Now, suppose that $q_i$ appears in $\mpath{w}$
    and it is 
    neither the initial state nor the accepting state.
    Let $r$ be the accepting run of $\cM$ on $w$.
    Since $\cM$ is acyclic, $q_i$ appears only once, which implies that $r$ is of the form
    \begin{equation*}
        \cdots p\transs{a} q_i \transs{a'}p'\cdots
    \end{equation*}
    Suppose that
    $\tuple{p, a, q_i}$ is the $j_1^{th}$ transition and 
    $\tuple{q_i, a', p'}$ is the $j_2^{th}$ transition.
    Then we have:
    \begin{equation*}
        \left(\bfA_\cM \cdot \mpath{w}\right)_i
        =
        \left(\bfA_\cM\right)_{i, j_1} +
        \left(\bfA_\cM\right)_{i, j_2}
        = 1 + (-1)
        = 0.
    \end{equation*}

    We will now show that
    $\kstar{\parikh{\cM}} = \sol{\cQ_\cM(\bfx)}$.

    \textbf{\uline{$\kstar{\parikh{\cM}} \subseteq \sol{\cQ_\cM(\bfx)}$}.}
    Suppose that $\bfv \in \kstar{\parikh{\cM}}$.
    By the definition of Kleene star and Parikh image,
    there exist accepted words $w_1, \ldots, w_\ell$
    and $\bfv = \sum_{i \in \intsinterval{\ell}} \parikh{w_i}$.
    Let $\bfu := \sum_{i \in \intsinterval{\ell}} \mpath{w_i}$.
    By the claim above and Lemma~\ref{lem:path}, we have
    \begin{equation*}
        \begin{aligned}
            \bfA_\cM \bfu\ =\ &\sum_{i \in \intsinterval{\ell}} \bfA \cdot \mpath{w_i}\ =\ \bfzero \\
            \bfB_\cM \bfu\ =\ &\sum_{i \in \intsinterval{\ell}} \bfB \cdot \mpath{w_i}
            \ =\ \sum_{i \in \intsinterval{\ell}} \parikh{w_i}\ =\ \bfv,
        \end{aligned}
    \end{equation*}
    which implies that concatenation of $\bfv$ and $\bfu$ is a solution of the integer linear system $\cQ_\cM(\bfx)$.
    Hence, $\bfv \in \sol{\cQ_\cM(\bfx)}$.

    \textbf{\uline{$\sol{\cQ_\cM(\bfx)}\subseteq\kstar{\parikh{\cM}}$}.}
    Suppose that $\bfv \in \sol{\cQ_\cM(\bfx)}$.
    Let $\bfu$ be the vector satisfying $\bfv = \bfB_\cM\bfu$ and $\bfA_\cM\bfu = \bfzero$.
    We interpret the $j^{th}$ entry of $\bfu$, denoted by $u_j$
    as a flow of the $j^{th}$ edge (transition) in the transition diagram of $\cM$.
    Then, since $\bfA_\cM\bfu = \bfzero$ by the claim above,
    for every node (state) $q$,
    if $q$ is neither the initial state nor the accepting state,
    then 
    \begin{equation*}
        \sum_{\substack{
                j \in \intsinterval{t}\\
                \text{the $j^{th}$ transition is of the form $(p, a, q)$}
        }} u_j
        -
        \sum_{\substack{
                j \in \intsinterval{t}\\
                \text{the $j^{th}$ transition is of the form $(q, a, p)$}
        }} u_j
        = 0.
    \end{equation*}
    That is, the incoming flow to the node (state) $q$ equals its outgoing flow.
    By the flow decomposition theorem~\cite{flow},
    there exist paths (words) $w_1, \ldots, w_\ell$ from the initial state to the accepting state
    such that $\bfu = \sum_{i \in \intsinterval{\ell}} \mpath{w_i}$.
    Since each $w_i$ is a path from the initial state to the accepting state, each
    $w_i$ is a word accepted by $\cM$.
    By Lemma~\ref{lem:path}, for each $i \leq \ell$, $\parikh{w_i} = \bfB_\cM\cdot\mpath{w_i}$.
    Thus, we have
    \begin{equation*}
        \bfv
        \ =\ \bfB_\cM\bfu
        \ =\ \sum_{i \in \intsinterval{\ell}} \bfB_\cM\cdot\mpath{w_i}
        \ =\ \sum_{i \in \intsinterval{\ell}} \parikh{w_i}
        \ \in\ \kstar{\parikh{\cM}}.
    \end{equation*}
\end{proof}

Combining Lemma~\ref{lem:ilp_to_nfa} and Lemma~\ref{lem:kstar_nfa} we obtain
a lemma regarding the Kleene star of the set of minimal solutions of an integer linear system.
\begin{lemma}\label{lemma:kleene_minimal}
    For every integer linear system $\cQ(\bfx): \bfA \bfx = \bfc$,
    where $\bfA \in \bbZ^{m \times n}$, $\bfc \in \bbZ^{m}$,
    there exists $\bfAtilde \in \bbZ^{(n+t) \times (n+k)}$ with $\norm{\bfAtilde} = 1$
    such that
    \begin{equation*}
        \kstar{\solmin{\cQ(\bfx)}}
        \ \subseteq\ 
        \sol{\exists \bfy\ (\bfx = \projp{n}{\bfy}) \land (\tilde{\bfA} \bfy = \bfzero)}
        \ \subseteq\ 
        \kstar{\sol{\cQ(\bfx)}},
    \end{equation*}
    where 
    $k = n\cdot(2D+1)^{m} + 1$,
    $t \le 2nk$, and
    $D = n \cdot \norm{\bfA} \cdot \left((n+1)\cdot \norm{\bfA} + \norm{\bfc} + 1\right)^{m}$.
    Moreover, $\bfAtilde$ can be computed in time $2^\bigO{\log n + \log K + m^2}$,
    where $K := \max\left(\norm{\bfA}, \norm{\bfc}\right)$.
\end{lemma}

We are now ready to prove Lemma \ref{lem:key}.
\begin{proof}
    Let $\cQ(\bfx): \bfA\bfx = \bfc$.
    Let $\bfAtilde$ be the matrix obtained by applying Lemma~\ref{lemma:kleene_minimal} to the integer linear system $\cQ(\bfx)$.
    It is routine to verify that the dimension and maximal absolute values of the entries of $\bfAtilde$ satisfy the required bounds, and $\bfAtilde$ can be computed within the desired time complexity.

    Here we show that $\kstar{\sol{\cQ(\bfx)}} = \sol{\cQtilde(\bfx)}$, where
    \begin{equation*}
        \cQtilde(\bfx)\ :=\ 
        \exists \bfy_1, \bfy_2\ 
        (\bfx\ =\ \projp{n}{\bfy_1} + \bfy_2)\ \land\ 
        (\bfAtilde\bfy_1\ =\ \bfzero)\ \land\ 
        (\bfA\bfy_2\ =\ \bfzero)\ \land\ 
        ((\bfy_1=\bfzero) \to (\bfy_2=\bfzero)).
    \end{equation*}

    \uline{$\kstar{\sol{\cQ(\bfx)}} \subseteq \sol{\cQtilde(\bfx)}$}.
    Suppose that $\bfv \in \kstar{\sol{\cQ(\bfx)}}$.
    If $\bfv = \bfzero$, then we can choose $\bfzero$
    as the assignment for $\bfy_1$ and $\bfy_1$.
    Otherwise if $\bfv \neq \bfzero$,
    by definition of Kleene star and minimal elements,
    there exists $\bfu_1, \bfu_2, \ldots, \bfu_k \in \sol{\cQ(\bfx)}$ and
    $\bfu_{1, \min}, \bfu_{2, \min}, \ldots, \bfu_{k, \min} \in \solmin{\cQ(\bfx)}$
    such that $\bfv = \sum_{i \in \intsinterval{k}} \bfu_i$,
    and for $1 \le i \le k$, $(\bfu_i - \bfu_{i, \min}) \in \bbN^n$.
    \begin{itemize}
        \item
            Because $\sum_{i \in \intsinterval{k}} \bfu_{i, \min} \in \kstar{\solmin{\cQ(\bfx)}}$,
            by Lemma~\ref{lemma:kleene_minimal},
            there exists
            $\bfv_1$ such that
            $\sum_{i \in \intsinterval{k}} \bfu_{i, \min} = \projp{n}{\bfv_1}$ and
            $\bfAtilde\bfv_1 = \bfzero$.
        \item
            Let $\bfv_2 = \sum_{i \in \intsinterval{k}} (\bfu_{i} - \bfu_{i, \min}$).
            Observe that
            $\bfA \bfv_2
            = \sum_{i \in \intsinterval{k}} \left(\bfA \bfu_{i} - \bfA \bfu_{i, \min}\right)
            = \bfzero$.
        \item
            Because $\bfc \neq \bfzero$, $\bfzero$ is not a solution of $\cQ(\bfx)$,
            we have $\bfv_1 \neq \bfzero$.
            Hence  $(\bfv_1=\bfzero) \to (\bfv_2=\bfzero)$ holds.
    \end{itemize}
    Since $\bfv = \projp{n}{\bfv_1} + \bfv_2$,
    $\bfv$ is a solution of $\cQtilde$, where the solution
    assigns $\bfy_1$ to $\bfv_1$ and assigns $\bfy_2$ to $\bfv_2$.

    \uline{$\sol{\cQtilde(\bfx)} \subseteq \kstar{\sol{\cQ(\bfx)}}$}.
    Suppose that $\bfv$ is a solution of $\cQtilde(\bfx)$
    by assigning $\bfv_1$  to $\bfy_1$ and assigning $\bfv_2$ to $\bfy_2$.
    \begin{itemize}
        \item
            If $\bfv_1 = \bfzero$, then $\bfv_2 = \bfzero$,
            which implies that
            $\bfv = \projp{n}{\bfv_1} + \bfv_2 = \bfzero$.
            By definition of Kleene star,
            $\bfzero \in \kstar{\sol{\cQ(\bfx)}}$.
        \item
            Otherwise if $\bfv_1 \neq \bfzero$,
            by Lemma~\ref{lemma:kleene_minimal},
            there exists $\bfu_1, \bfu_2, \ldots, \bfu_k \in \sol{\cQ(\bfx)}$ such that
            $\projp{n}{\bfv_1} = \sum_{i \in \intsinterval{k}} \bfu_i$.
            Note that because $\bfA\bfv_2 = \bfzero$, we have
            \begin{equation*}
                \bfA (\bfu_1 + \bfv_2)
                \ =\ \bfA\bfu_1 + \bfA\bfv_2
                \ =\ \bfc + \bfzero
                \ =\ \bfc,
            \end{equation*}
            which implies that $\bfu_1 + \bfv_2$ is a solution of $\cQ(\bfx)$.
            Therefore, we infer
            \begin{equation*}
                \bfv_1 + \bfv_2
                \ =\ (\bfu_1 + \bfv_2) + \sum_{i \in \intinterval{2}{k}} \bfu_i
                \ \in\ \kstar{\sol{\cQ(\bfx)}}.
            \end{equation*}
            The membership in $\kstar{\sol{\cQ(\bfx)}}$
            follows from the fact that $\bfu_1 + \bfv_2,\bfu_2,\ldots,\bfu_k$ are all solutions of $\cQ$.
    \end{itemize}
\end{proof}
